% Naskah utama — makalah revisi.
% Kompilasi: lualatex + biber + lualatex + lualatex.
\documentclass[12pt,a4paper]{article}

\usepackage{fontspec}
\setmainfont{TeX Gyre Termes}
\usepackage{geometry}
\geometry{a4paper, top=2.54cm, bottom=2.54cm, left=2.54cm, right=2.54cm}
\setlength{\headheight}{14pt}
\usepackage{setspace}
\singlespacing
\usepackage{booktabs}
\usepackage{tabularx}
\usepackage{array}
\newcolumntype{Y}{>{\raggedright\arraybackslash}X}
\usepackage{enumitem}
\usepackage{caption}
\captionsetup{font=small, labelfont=bf, skip=6pt}
\usepackage{microtype}
\usepackage{fancyhdr}
\usepackage{csquotes}
\usepackage[notes,backend=biber]{biblatex-chicago}
\addbibresource{referensi.bib}
\usepackage{xurl}
\usepackage[hidelinks]{hyperref}
\usepackage{graphicx}
\renewcommand{\contentsname}{Daftar Isi}
\renewcommand{\listtablename}{Daftar Tabel}
\renewcommand{\tablename}{Tabel}
\setlength{\emergencystretch}{3em}

\pagestyle{fancy}
\fancyhf{}
\fancyhead[L]{\small\itshape Pengorganisasian Pengetahuan}
\fancyhead[R]{\small\itshape al-Farabi --- Basis Data}
\fancyfoot[C]{\small\thepage}
\renewcommand{\headrulewidth}{0.4pt}
\fancypagestyle{plain}{%
  \fancyhf{}%
  \fancyfoot[C]{\small\thepage}%
  \renewcommand{\headrulewidth}{0pt}%
}

\begin{document}

% ===== Halaman sampul =====
\begin{titlepage}
\centering
{\large\bfseries UNIVERSITAS DARUSSALAM GONTOR\par}
{\normalsize Fakultas Sains dan Teknologi\\Program Studi Teknik Informatika\par}
\vspace{1.4cm}
{\rule{\textwidth}{1pt}\par}
\vspace{0.7cm}
{\Large\bfseries Pengorganisasian Pengetahuan dalam Peradaban Islam\\dan Basis Data Modern:\par}
\vspace{0.25cm}
{\Large Perbandingan Konseptual melalui Logika al-Farabi\par}
\vspace{0.7cm}
{\rule{\textwidth}{1pt}\par}
\vspace{0.9cm}
{\normalsize Makalah Penelitian\par}
\vspace{0.3cm}
{\small Fokus: kontribusi Abu Nasr al-Farabi dalam logika serta relevansinya\\dengan logika formal dalam Teknik Informatika\par}
\vspace{1.4cm}
{\normalsize Disusun oleh:\par}
\vspace{0.2cm}
{\large\bfseries Farrel Ghozy Affifudin\par}
{\normalsize NIM 452024611053\par}
\vfill
{\normalsize 2026\par}
\end{titlepage}

% ===== Abstrak =====
\section*{Abstrak}
\addcontentsline{toc}{section}{Abstrak}
Makalah ini membahas pengorganisasian pengetahuan dalam peradaban Islam dan basis data modern melalui perbandingan konseptual yang bertumpu pada logika Abu Nasr al-Farabi. Masalah yang diangkat adalah bagaimana tradisi perpustakaan dan pengatalogan manuskrip berkembang, apa kontribusi spesifik al-Farabi dalam logika, serta apa persamaan dan batas perbandingannya dengan logika formal dalam Teknik Informatika. Metode yang digunakan adalah kajian pustaka dengan pendekatan perbandingan konseptual. Sumber utama adalah kajian logika al-Farabi dan logika Aristoteles, kajian filsafat informasi, serta kajian model data relasional, dilengkapi literatur sejarah perpustakaan Islam sebagai konteks. Hasil analisis menunjukkan bahwa katalog manuskrip dan katalog sistem basis data memiliki kemiripan fungsional sebagai sarana temu kembali, tetapi berbeda fondasi teknis. Logika al-Farabi, khususnya teori silogisme, demonstrasi, dan dialektika, menyediakan jembatan konseptual menuju representasi dan inferensi dalam sistem komputasi, tetapi silogisme dalam bahasa alamiah yang didisiplinkan tidak sama dengan kalkulus predikat yang dapat dieksekusi mesin. Simpulan makalah ini adalah perbandingan tersebut sah sebagai analisis konseptual dan gugur sebagai klaim hubungan sejarah langsung.

\noindent\textbf{Kata kunci:} pengorganisasian pengetahuan, al-Farabi, logika Arab, katalog manuskrip, basis data relasional, temu kembali informasi.

\newpage
\tableofcontents
\listoftables
\newpage

\section{Pendahuluan}

Perpustakaan dalam peradaban Islam bukan sekadar tempat penyimpanan buku. Survei Algeriani dan Mohadi mencatat bahwa Bayt al-Hikmah di Baghdad memiliki sistem manajemen yang sangat terorganisasi, khususnya dalam pengumpulan dan pengatalogan buku.\autocite{algeriani-mohadi2017} Pernyataan tersebut penting karena menunjukkan bahwa masalah pengorganisasian pengetahuan, yaitu bagaimana menemukan kembali apa yang telah disimpan, sudah dilembagakan jauh sebelum basis data modern. Perhatian kajian modern terhadap perpustakaan Islam pun cukup tua: kajian tentang empat perpustakaan besar Baghdad abad pertengahan telah terbit sejak 1932.\autocite{mackensen1932}

Di sisi lain, Abu Nasr al-Farabi, filsuf abad ke-10 M, menempati kedudukan sentral dalam pewarisan dan pengembangan logika Aristoteles dalam bahasa Arab.\autocite{druart2024} Kajian Hodges dan Druart menunjukkan bahwa komentar-komentarnya atas \emph{Organon} serta uraian sistematisnya tentang silogisme, demonstrasi, dan dialektika menjadikannya titik rujuk utama tradisi logika Arab-Islam.\autocite{hodges-druart2019} Kedudukan edisi kritis komentarnya atas \emph{De Interpretatione} yang disunting Zimmermann pada 1981 diketahui melalui kajian tersebut.\autocite{zimmermann1981} Kajian modern menempatkan silogisme Arab, yang mencakup silogisme kategoris, modal, dan hipotetik, sebagai bidang tersendiri dari al-Farabi hingga Averroes.\autocite{chatti2019} Tradisi ini bahkan memperluas cakupan \emph{Organon} hingga mencakup retorika dan puisi.\autocite{black1990} Karya dialektika al-Farabi, \emph{Kitab al-Jadal}, telah dikaji dalam monograf tersendiri,\autocite{dipasquale2019} dan terjemahan Inggris atas karya-karya filsafatnya telah tersedia sejak 1969.\autocite{mahdi1969}

Relevansinya bagi Teknik Informatika terletak pada kedudukan logika formal. Model data relasional dibangun di atas kalkulus predikat terapan,\autocite{codd1970} temu kembali informasi telah menjadi disiplin buku teks tersendiri,\autocite{manning2008} dan filsafat informasi kini diakui sebagai subdisiplin filsafat.\autocite{adriaans2023} Makalah ini menjawab empat pertanyaan: (1) bagaimana tradisi perpustakaan dan pengatalogan berkembang dalam peradaban Islam; (2) apa kontribusi spesifik al-Farabi terhadap logika; (3) apa persamaan dan perbedaan logika al-Farabi dengan logika formal dalam Teknik Informatika; dan (4) sejauh mana perbandingan itu sah dan di mana batasnya.

Tesis makalah ini adalah kemiripan antara pengorganisasian pengetahuan klasik dan modern bersifat fungsional, yaitu menjawab masalah temu kembali yang sejenis, bukan genealogis dalam arti garis keturunan sejarah langsung. Logika menjadi jembatan konseptual yang sah, tetapi perbedaan fondasi membatasi analogi tersebut secara tegas.

\section{Metode Penelitian}

Makalah ini menggunakan kajian pustaka dengan pendekatan perbandingan konseptual. Langkah kerjanya adalah menghimpun literatur tentang perpustakaan Islam, logika al-Farabi, dan logika formal dalam komputasi, kemudian membandingkan konsep-konsepnya berdasarkan fungsi dan fondasinya, bukan berdasarkan dugaan pengaruh sejarah.

Cakupan sumber terdiri atas dua kelompok. Kelompok pertama adalah sumber yang diperiksa teks penuhnya: entri \emph{Stanford Encyclopedia of Philosophy} tentang logika dan bahasa al-Farabi, logika Aristoteles, Ibnu Sina, dan informasi, serta dua artikel jurnal tentang Bayt al-Hikmah dan model data relasional. Kelompok kedua adalah literatur pendukung yang digunakan untuk konteks umum tanpa kutipan halaman karena teks penuhnya belum diakses: kajian tentang silogisme Arab, retorika dan puisi dalam filsafat Arab, dialektika al-Farabi, terjemahan karya al-Farabi, gerakan penerjemahan Graeco-Arab, budaya buku Arab, kajian awal perpustakaan Baghdad, dan buku teks temu kembali informasi.

Naskah ini membedakan tiga hal: fakta yang didukung referensi, interpretasi penulis yang dinyatakan eksplisit sebagai interpretasi penulis, dan analogi konseptual yang dinyatakan eksplisit sebagai analogi. Keterbatasannya adalah klaim dari kelompok kedua hanya bersifat umum. Angka-angka rinci tentang tahun pendirian atau jumlah koleksi perpustakaan tidak dibahas karena memerlukan silang dengan sumber primer seperti \emph{Fihrist} Ibn al-Nadim dan teks penuh kajian sejarah yang belum diakses.

\section{Kajian Historis: Perpustakaan dan Pengetahuan dalam Peradaban Islam}

\subsection{Bayt al-Hikmah dan tradisi pengatalogan}

Konteks intelektual Baghdad abad ke-8 hingga ke-10 M dibentuk oleh gerakan penerjemahan karya Yunani ke bahasa Arab.\autocite{gutas1998} Dalam ekosistem inilah institusi ilmu seperti Bayt al-Hikmah tumbuh. Menurut survei Algeriani dan Mohadi, institusi ini memiliki sistem manajemen yang sangat terorganisasi, terutama dalam pengumpulan dan pengatalogan buku.\autocite{algeriani-mohadi2017}

Batas sumber ini perlu dinyatakan. Artikel tersebut merupakan survei sekunder sehingga makalah ini hanya mengambil klaim umum tentang keberadaan praktik pengatalogan yang sistematis. Budaya buku Arab secara umum, sebagai wadah tradisi manuskrip tersebut, dibahas dalam kajian Pedersen sebagai konteks.\autocite{pedersen1984}

\subsection{al-Farabi dalam tradisi logika}

Al-Farabi hidup dalam dunia yang dibentuk gerakan penerjemahan tersebut. Korpus logika Aristoteles telah tersedia dalam bahasa Arab dan membutuhkan sistematisasi. Gambaran umum riwayat intelektualnya merujuk pada entri umum Druart, sedangkan seluruh klaim logika yang spesifik merujuk pada kajian Hodges dan Druart.\autocite{druart2024}\autocite{hodges-druart2019}

Kedudukan al-Farabi bukan sekadar komentator pasif. Pertama, keberadaan edisi kritis atas komentarnya terhadap \emph{De Interpretatione} menunjukkan bobot filologis karya logikanya.\autocite{zimmermann1981} Kedua, kajian kontemporer menempatkan kontribusinya dalam busur panjang silogisme Arab dari al-Farabi hingga Averroes yang mencakup silogisme kategoris, modal, dan hipotetik.\autocite{chatti2019} Ketiga, kurikulum logika sembilan bagian, mulai dari \emph{Eisagoge}, \emph{Categories}, \emph{De interpretatione}, \emph{Syllogism}, \emph{Demonstration}, \emph{Dialectic}, \emph{Sophistics}, hingga \emph{Rhetoric} dan \emph{Poetics}, berlanjut ke Ibnu Sina. Kelanjutan ini menandakan kesinambungan tradisi yang dirintis generasi al-Farabi, bukan kontribusi yang lewat sekali.\autocite{gutas2025}

\section{Analisis Kontribusi Keilmuan: Logika al-Farabi}

\subsection{Logika sebagai disiplin penilaian penalaran}

Al-Farabi menempatkan logika sebagai disiplin yang menilai penalaran, sejajar dengan kedudukan tata bahasa bagi bahasa. Dalam \emph{Ihsa al-Ulum} ia menyatakan bahwa subjek logika adalah konsep-konsep sejauh diungkapkan oleh ungkapan, dan ungkapan-ungkapan sejauh mengungkapkan konsep.\autocite[bag.~5]{hodges-druart2019} Definisi ini penting bagi argumen makalah: sejak awal logika al-Farabi sudah mempertemukan konsep, ungkapan, dan aturan. Sebagai interpretasi penulis, tiga kutub tersebut dapat dianalogikan dengan representasi, sintaks, dan inferensi dalam sistem komputasi, tetapi analogi ini dibahas tersendiri pada bagian perbandingan.

Ia juga memberi definisi silogisme yang menekankan tatanan mental: silogisme adalah hal-hal yang tersusun dalam pikiran dalam suatu tatanan.\autocite[bag.~7]{hodges-druart2019}

\subsection{Silogisme kategoris, hipotetik, dan demonstrasi}

Latarnya adalah doktrin Aristoteles tentang deduksi, yaitu ujaran yang di dalamnya, dengan mengandaikan hal-hal tertentu, dihasilkan sesuatu yang berbeda secara niscaya karena hal-hal tersebut demikian.\autocite[bag.~3]{smith2022} Al-Farabi menerima kerangka ini lalu mengembangkannya. Ia mengakui tiga figur silogisme kategoris dengan empat, empat, dan enam modus yang valid,\autocite[bag.~6.1]{hodges-druart2019} dan memberi contoh silogisme yang valid: tidak ada batu yang merupakan hewan, sebagian benda adalah hewan, maka sebagian benda bukan batu.\autocite[\emph{Syllogism} 27.6f, dikutip bag.~5]{hodges-druart2019}

Kontribusinya melampaui pengulangan. Pertama, ia memformalkan argumen majemuk hipotetik.\autocite[bag.~6.2]{hodges-druart2019} Kedua, ia menguraikan demonstrasi sebagai silogisme \emph{Barbara} yang diberi tambahan-tambahan.\autocite[bag.~6.4]{hodges-druart2019} Ketiga, menurut kajian Hodges dan Druart, ia tercatat paling awal secara eksplisit menyimpulkan muatan eksistensial dari pernyataan bentuk setiap B adalah A. Klaim keaslian ini dikutip sebagai atribusi kajian tersebut, bukan sebagai verifikasi independen penulis.\autocite[bag.~6.1]{hodges-druart2019}

Dimensi dialektika, yaitu seni perdebatan dan pengujian premis, diwakili oleh \emph{Kitab al-Jadal} yang kini tersedia kajian monografisnya.\autocite{dipasquale2019} Naskah ini tidak menyamakan dialektika tersebut dengan penalaran \emph{defeasible} modern. Kesetaraan seperti itu memerlukan pembahasan tersendiri dan dukungan sumber yang belum diakses dalam penelitian ini.

\subsection{Keberatan: apakah kontribusi al-Farabi dilebih-lebihkan?}

Keberatan yang wajar adalah penonjolan al-Farabi berisiko mengecilkan Ibnu Sina dan Ibnu Rusyd. Jawabannya bersifat proporsional. Makalah ini tidak menempatkan al-Farabi sebagai puncak, melainkan sebagai simpul sistematisasi awal yang terdokumentasi. Justru kesinambungan kurikulum sembilan bagiannya ke Ibnu Sina menjadi bukti pengaruhnya yang terlembaga.\autocite{gutas2025} Kajian Chatti yang membentang hingga Averroes meneguhkan bingkai kesinambungan ini, bukan bingkai pemujaan.\autocite{chatti2019}

\section{Relevansi terhadap Teknik Informatika}

\subsection{Logika formal sebagai fondasi komputasi}

Fondasi yang dipakai Teknik Informatika berbeda secara kualitatif. Codd membangun model data relasional di atas kalkulus predikat terapan. Ia menyatakan bahwa adopsi model relasional memungkinkan pengembangan bahasa bagian data yang universal yang berlandaskan kalkulus predikat terapan.\autocite[381]{codd1970} Ia menuntut agar kelas ekspresi kualifikasi memiliki daya deskriptif seperti kelas formula berkuantor yang baik dalam kalkulus predikat terapan.\autocite[381]{codd1970} Relevan dengan tema perpustakaan, Codd menetapkan bahwa relasi-relasi yang dideklarasikan dimasukkan ke dalam katalog sistem.\autocite[381]{codd1970} Di sinilah titik temu fungsional yang presisi: katalog, yaitu daftar koleksi agar temu kembali dimungkinkan, muncul kembali sebagai komponen arsitektur basis data. Temu kembali informasi sendiri telah menjadi disiplin buku teks yang mapan,\autocite{manning2008} dan filsafat informasi, termasuk informasi semantik, diakui sebagai subdisiplin filsafat yang berhubungan erat dengan filsafat logika.\autocite{adriaans2023}

\subsection{Perbandingan terstruktur}

Tabel~\ref{tab:banding} merangkum perbandingan. Kolom terakhir menandai status setiap baris. Kode H berarti hubungan sejarah langsung, dan tidak ada baris yang memenuhi syarat bukti untuk kode ini. Kode F berarti kemiripan fungsional, sedangkan kode T berarti perbedaan teknis.

\begin{table}[htbp]
\centering
\small
\caption{Perbandingan logika al-Farabi dan logika formal dalam Teknik Informatika}
\label{tab:banding}
\renewcommand{\arraystretch}{1.2}
\begin{tabularx}{\textwidth}{@{}>{\raggedright\arraybackslash}p{2.4cm}>{\raggedright\arraybackslash}X>{\raggedright\arraybackslash}Xc@{}}
\toprule
Aspek & al-Farabi & Teknik Informatika & Status\\
\midrule
Objek penalaran & Menilai kebenaran argumen dalam bahasa alamiah & Menilai validitas komputasi dan kueri formal & F\\
Konsep dan definisi & Konsep diungkapkan ungkapan; klasifikasi ilmu & Ontologi, taksonomi, skema metadata & F\\
Bentuk inferensi & Silogisme kategoris dan hipotetik; demonstrasi \emph{Barbara} & Kalkulus proposisional dan predikat; resolusi & T\\
Katalog & \emph{Fihrist} sebagai alat bantu manusia menemukan manuskrip & Katalog sistem sebagai metadata yang dibaca mesin untuk optimasi kueri & F+T\\
Temu kembali & Klasifikasi dan daftar koleksi & Model Boolean dan probabilistik; indeks terbalik & F+T\\
Basis formal & Bahasa dan tatanan mental & Teori himpunan, aljabar relasional, dan indeks & T\\
\bottomrule
\end{tabularx}
\\[4pt]
{\footnotesize Status: H = hubungan sejarah langsung (tidak ada yang terbukti);
F = kemiripan fungsional; T = perbedaan teknis. Tabel merupakan sintesis penulis dari kajian logika, kajian perpustakaan, dan kajian basis data yang dirujuk makalah ini.}
\end{table}

Dua baris perlu diuraikan. Pertama, katalog. \emph{Fihrist} manuskrip dan katalog sistem Codd sama-sama menyimpan data tentang data agar koleksi dapat ditemukan kembali, sehingga ada kemiripan fungsional. Namun katalog Codd dibaca oleh mesin pengoptimasi kueri dan terikat pada kalkulus formal, sedangkan \emph{fihrist} dibaca manusia dan terikat konvensi filologis. Itu adalah perbedaan teknis yang tidak boleh dikaburkan. Kedua, demonstrasi \emph{Barbara} dan pembuktian teorema. Keduanya mengejar kepastian dari premis, tetapi demonstrasi al-Farabi bergerak dalam bahasa alamiah yang didisiplinkan, sedangkan pembuktian komputasi bergerak dalam bahasa formal yang dapat dieksekusi mesin.

\subsection{Batas perbandingan}

Perbandingan di atas memiliki tiga batas. Pertama, tidak ada hubungan sejarah langsung. Tidak ada bukti bahwa al-Farabi memengaruhi Codd atau perancang basis data. Kesamaan kosakata katalog adalah konvergensi fungsional, dan mengklaim sebaliknya adalah anakronisme. Kedua, klasifikasi ilmu pra-modern bersifat normatif dan filosofis karena menetapkan tatanan ilmu yang seharusnya, sedangkan skema basis data bersifat operasional karena harus mengeksekusi kueri dalam batas waktu. Ketiga, basis data tidak dapat direduksi menjadi logika semata. Model data, struktur indeks, bahasa kueri, dan rancangan sistem adalah lapisan teknik yang tidak identik dengan silogisme. Arsitektur Codd sendiri memisahkan bahasa bagian data dari kalkulusnya.\autocite[381]{codd1970}

\section{Kesimpulan}

Menjawab empat rumusan masalah, makalah ini menyimpulkan empat hal. Pertama, peradaban Islam melembagakan pengorganisasian pengetahuan melalui institusi seperti Bayt al-Hikmah dengan praktik pengatalogan yang sistematis. Kedua, kontribusi spesifik al-Farabi terletak pada definisi subjek logika, sistematisasi tiga figur silogisme, formalisasi argumen hipotetik majemuk, teori demonstrasi, dan perhatian pada dialektika. Ketiga, persamaannya dengan logika dalam Teknik Informatika bersifat fungsional, yaitu kesamaan masalah temu kembali dan penilaian penalaran, sedangkan perbedaannya bersifat fondasional, yaitu silogisme dalam bahasa alamiah berbanding kalkulus formal yang dieksekusi mesin. Keempat, perbandingan ini sah sebagai analisis konseptual dan gugur sebagai klaim sejarah.

Keterbatasan yang masih terbuka adalah perlunya teks penuh kajian sejarah dan kajian logika Arab untuk menguatkan klaim yang saat ini baru didukung sebagai konteks umum, serta perlunya kajian atas \emph{Fihrist} Ibn al-Nadim untuk rincian praktik pengatalogan klasik.

\clearpage
\printbibliography[title=Daftar Pustaka]

\end{document}
